\appendix
\section{Partial-wave projection of one-meson exchange}
\label{app:ome-projection}

We project the time-ordered one-meson-exchange potentials with the spin-angular
projectors, using the spin average and the two normalised solid-angle measures
of Ref.~\cite{Zhang:2026iyr}, v1, Eq.~(80). After contracting all polarisation
indices, rotational invariance reduces the double angular average to
$\tfrac12\int_{-1}^{1}dx$, with $x=\hat{\bm p}'\cdot\hat{\bm p}$.
The remaining numerator is a polynomial in $x$, so each entry is a finite
linear combination of moments
\begin{equation}
 I_k^{(\beta)}(p',p;E,m_\phi)
 =\int_{-1}^{1}\frac{x^k\,dx}{2\omega(\omega+\nu_\beta)},
 \qquad
 \omega=\sqrt{p^2+p'^2+m_\phi^2-2pp'x}.
 \label{eq:ome-angular-moments}
\end{equation}
Here $\nu_\beta$ denotes the intermediate two-heavy-meson energy minus the
absolute energy $M$; it is the quantity called $\mu_\beta$ in
Ref.~\cite{Zhang:2026iyr}, v1, Eqs.~(76)--(77), not a two-body reduced mass.
The heavy-meson energies there are written as
$\omega_{B^{(*)}}(p)=m_{B^{(*)}}+p^2/(2m_{B^{(*)}})$ here, to distinguish
them from the binding energy $E_B$.
Retaining the mass splitting, recoil and energy dependence, this procedure
reproduces all nine $1^{++}$ entries of that reference, v1,
Eqs.~(D2)--(D10), with exactly the projector phases of its Eq.~(D1).
As an independent check, direct numerical integration of both solid angles
in Eq.~(80), with the three-dimensional spin matrices, agrees entry by entry
with the analytic kernel for pion and eta exchange separately at three seeded
random points below the three-body thresholds. The largest absolute residual
is $4.7\times10^{-13}\,\GeV^{-2}$; increasing the angular order changes the
result by at most $1.5\times10^{-13}\,\GeV^{-2}$ at these points.
The isoscalar combination is $3V(m_\pi)+V(m_\eta)/3$ as in that reference,
v1, Eq.~(82).

For $J^{PC}=2^{++}$ the six channels are, in order,
$B^*\bar B^*({}^5S_2)$, $B^*\bar B^*({}^1D_2)$, $B^*\bar B^*({}^5D_2)$,
$B^*\bar B^*({}^5G_2)$, $B\bar B^*({}^3D_2)$ and $B\bar B({}^1D_2)$, with the
$C$-even $B\bar B^*$ combination
$(|B\bar B^*\rangle-|B^*\bar B\rangle)/\sqrt2$ and the Cartesian projectors of
Ref.~\cite{Baru:2019xnh}, v2, App.~B, with their printed signs. The energy is
measured from $2m_{B^*}$, $M=2m_{B^*}+E$, so the offsets of
Ref.~\cite{Zhang:2026iyr}, v1, Eq.~(77) become
$\nu_1=-2\Delta m+(p^2+p'^2)/(2m_B)-E$,
$\nu_2=(p^2+p'^2)/(2m_{B^*})-E$,
$\nu_3=-\Delta m+p^2/(2m_B)+p'^2/(2m_{B^*})-E$ and
$\nu_4=-\Delta m+p'^2/(2m_B)+p^2/(2m_{B^*})-E$, with $\Delta m=m_{B^*}-m_B$.
Besides the potentials of Ref.~\cite{Zhang:2026iyr}, v1, Eqs.~(72)--(75), the
Lagrangian of its Eq.~(68) gives the $B\bar B\leftrightarrow B^*\bar B^*$
transition
\begin{equation}
 V^{\rm(OME)}_{B\bar B\to B^*\bar B^*}(\bm p',\bm p;m_\phi)
 =\frac{g_b^2}{4F_\pi^2}(\bm\epsilon_1'^*\cdot\bm q)(\bm\epsilon_2'^*\cdot\bm q)\,
 2D_3(\bm p',\bm p;m_\phi),
 \label{eq:ome-bb-bsbs}
\end{equation}
and the reverse process with $\bm\epsilon_{1,2}$ and $2D_4$. Both time
orderings have the same intermediate pair $B(p)\bar B^*(p')$ or $B^*(p')\bar
B(p)$, hence the factor two. The sign is fixed relative to Eq.~(72) of that
reference: the meson-line vertices $B\to B^*$ and $B^*\to B$ of Eq.~(68) carry
the same coefficient, and with the antimeson line in common the $C$-even
$B\bar B^*$ potential equals minus the corresponding cross amplitude. There is
no $B\bar B\to B\bar B$ or $B\bar B\to B\bar B^*$ one-boson exchange. Each
projected entry is
\begin{equation}
 V^{\rm(OME)}_{\alpha'\alpha}(p',p;m_\phi)
 =-\frac{g_b^2}{8F_\pi^2}\sum_{k=0}^{4}c_{\alpha'\alpha,k}(p',p)
 \bigl[I_k^{(\beta)}+I_k^{(\gamma)}\bigr],
 \label{eq:ome-2pp-entry}
\end{equation}
with $(\beta\gamma)=(22)$ within $B^*\bar B^*$, $(12)$ within $B\bar B^*$,
$(23)$ from $B\bar B^*$ to $B^*\bar B^*$, $(24)$ for the reverse, and $(33)$
from $B\bar B$ to $B^*\bar B^*$, $(44)$ for the reverse. The numerator
polynomials of the upper triangle are listed in Table~\ref{tab:ome-2pp}; the
lower triangle follows from $c_{\alpha\alpha'}(p,p')=c_{\alpha'\alpha}(p',p)$.
The $G$-wave entries need moments up to $x^4$, one more than in $1^{++}$.
The entries ${}^5S_2$--${}^5G_2$ and ${}^1D_2$--$B\bar B^*({}^3D_2)$ vanish
identically, as do the two $B\bar B$ entries without a $B^*\bar B^*$ leg.

\begin{table}[tbp]
\centering
\caption{Numerator coefficients $c_{\alpha'\alpha,k}$ of
Eq.~\eqref{eq:ome-2pp-entry} for $2^{++}$, with $r=p'$ outgoing and $p$
incoming. Channels $1$--$6$ are $B^*\bar B^*({}^5S_2,{}^1D_2,{}^5D_2,{}^5G_2)$,
$B\bar B^*({}^3D_2)$ and $B\bar B({}^1D_2)$. With
$s_k=\bigl(-\tfrac{2p^2-r^2}{6},\tfrac{2pr}{3},-\tfrac{r^2}{2}\bigr)_k$, the
entries $12,13,15,16$ are $s_k$ times $\sqrt{2/5}$, $-\sqrt{14/5}$,
$-\sqrt{18/5}$ and $-\sqrt{6/5}$; entries $14,25,56,66$ vanish.}
\label{tab:ome-2pp}
\small
\begin{tabular}{@{}cccccc@{}}
\toprule
$\alpha'\alpha$ & $c_0$ & $c_1$ & $c_2$ & $c_3$ & $c_4$\\
\midrule
11 & $\frac{p^2+r^2}{3}$ & $-\frac{2pr}{3}$ & $0$ & $0$ & $0$\\
22 & $\frac{p^2+r^2}{3}$ & $-\frac{2pr}{3}$ & $-(p^2+r^2)$ & $2pr$ & $0$\\
23 & $-\frac{\sqrt7(p^2+r^2)}{21}$ & $-\frac{\sqrt7pr}{21}$ & $\frac{\sqrt7(p^2+r^2)}{7}$ & $-\frac{\sqrt7pr}{7}$ & $0$\\
24 & $\frac{\sqrt{35}(4p^2-3r^2)}{140}$ & $-\frac{6\sqrt{35}pr}{35}$ & $-\frac{3\sqrt{35}(2p^2-5r^2)}{70}$ & $\frac{2\sqrt{35}pr}{7}$ & $-\frac{\sqrt{35}r^2}{4}$\\
26 & $\frac{\sqrt3(p^2+r^2)}{6}$ & $-\frac{\sqrt3pr}{3}$ & $-\frac{\sqrt3(p^2+r^2)}{2}$ & $\sqrt3pr$ & $0$\\
33 & $-\frac{5(p^2+r^2)}{21}$ & $\frac{11pr}{42}$ & $\frac{5(p^2+r^2)}{7}$ & $-\frac{17pr}{14}$ & $0$\\
34 & $-\frac{\sqrt5(4p^2-3r^2)}{70}$ & $\frac{12\sqrt5pr}{35}$ & $\frac{3\sqrt5(2p^2-5r^2)}{35}$ & $-\frac{4\sqrt5pr}{7}$ & $\frac{\sqrt5r^2}{2}$\\
35 & $\frac{\sqrt7(p^2+r^2)}{14}$ & $\frac{\sqrt7pr}{14}$ & $-\frac{3\sqrt7(p^2+r^2)}{14}$ & $\frac{3\sqrt7pr}{14}$ & $0$\\
36 & $\frac{\sqrt{21}(p^2+r^2)}{21}$ & $\frac{\sqrt{21}pr}{21}$ & $-\frac{\sqrt{21}(p^2+r^2)}{7}$ & $\frac{\sqrt{21}pr}{7}$ & $0$\\
44 & $-\frac{3(p^2+r^2)}{56}$ & $-\frac{3pr}{7}$ & $\frac{15(p^2+r^2)}{28}$ & $\frac{5pr}{7}$ & $-\frac{5(p^2+r^2)}{8}$\\
45 & $-\frac{\sqrt{35}(3p^2-4r^2)}{70}$ & $-\frac{12\sqrt{35}pr}{35}$ & $\frac{3\sqrt{35}(5p^2-2r^2)}{35}$ & $\frac{4\sqrt{35}pr}{7}$ & $-\frac{\sqrt{35}p^2}{2}$\\
46 & $\frac{\sqrt{105}(3p^2-4r^2)}{140}$ & $\frac{6\sqrt{105}pr}{35}$ & $-\frac{3\sqrt{105}(5p^2-2r^2)}{70}$ & $-\frac{2\sqrt{105}pr}{7}$ & $\frac{\sqrt{105}p^2}{4}$\\
55 & $0$ & $\frac{pr}{2}$ & $0$ & $-\frac{pr}{2}$ & $0$\\
\bottomrule
\end{tabular}
\end{table}

Three checks test these $2^{++}$ expressions in the implemented kernel.
First, direct numerical integration of both solid angles in
Ref.~\cite{Zhang:2026iyr}, v1, Eq.~(80), with the Cartesian projectors and the
three-dimensional spin matrices including Eq.~\eqref{eq:ome-bb-bsbs}, agrees
with all 21 independent entries for pion and eta exchange separately at three
seeded off-shell points below the $B\bar B\pi$ threshold; the largest residual
is $2.9\times10^{-12}\,\GeV^{-2}$ and raising the angular order changes the
integral by at most $2.7\times10^{-15}\,\GeV^{-2}$. Second, time reversal,
$V_{\alpha'\alpha}(p',p;E)=V_{\alpha\alpha'}(p,p';E)$, holds to
$7\times10^{-16}\,\GeV^{-2}$; giving one $B^*\bar B^*\leftarrow B\bar B$ entry the
offset $\nu_4$ instead of $\nu_3$ breaks it by up to
$3\times10^{-4}\,\GeV^{-2}$. Third, in the strict heavy-quark spin symmetry
limit, $m_B=m_{B^*}=(m_B+3m_{B^*})/4$, the potential rotated by 9j and 6j
recoupling to heavy-quark spin $s_H$ and light angular momentum $j_\ell$
(including the orbital one) is block diagonal in $(s_H,j_\ell)$ to
$1.1\times10^{-15}\,\GeV^{-2}$, and its $(s_H,j_\ell)=(1,1)$ block equals the
corresponding $1^{++}$ block to $10^{-13}\,\GeV^{-2}$, as heavy-quark spin
symmetry requires. The rotation is checked against the recoupling coefficients
of Ref.~\cite{Bruschini:2023zkb}, v3, Table~2, and does not use the projectors;
reversing the sign of the $B\bar B({}^1D_2)$ projector destroys the block
structure.
